\documentclass[11pt]{article}
\usepackage[T1]{fontenc}
\usepackage[utf8]{inputenc}
\usepackage{lmodern,amsmath,amssymb,amsthm,mathtools}
\usepackage[margin=1.05in]{geometry}
\usepackage{microtype,enumitem,tikz-cd}
\usepackage[colorlinks=true,linkcolor=blue,citecolor=blue,urlcolor=blue]{hyperref}
\usepackage{xurl}
\setlength{\emergencystretch}{2em}
\newtheorem{theorem}{Theorem}[section]
\newtheorem{lemma}[theorem]{Lemma}
\newtheorem{proposition}[theorem]{Proposition}
\newtheorem{corollary}[theorem]{Corollary}
\theoremstyle{remark}
\newtheorem{remark}[theorem]{Remark}
\newcommand{\C}{\mathbb C}
\newcommand{\OO}{\mathcal O}
\DeclareMathOperator{\ord}{ord}
\DeclareMathOperator{\Supp}{Supp}
\numberwithin{equation}{section}
\title{Anticanonical nonvanishing from smooth semipositivity}
\author{OpenAI}
\date{September 26, 2026}
\hypersetup{pdftitle={Anticanonical nonvanishing from smooth semipositivity},pdfauthor={OpenAI},pdfsubject={Anticanonical-nonvanishing-from-smooth-semipositivity-September-26-2026}}
\begin{document}
\maketitle
\begin{abstract}
Every smooth connected projective complex variety with smoothly
semipositive anticanonical bundle has a nonzero section of some positive
anticanonical power.
\end{abstract}
\tableofcontents
\section{Introduction}
\label{sec:introduction}

For a smooth projective complex variety \(X\), the anticanonical line
bundle is \(-K_X=\det T_X\). A smooth Hermitian metric on this line is
semipositive when its Chern curvature is a nonnegative real \((1,1)\)-form.
Such a metric makes the numerical class nef. The nonvanishing question
asks whether \(H^0(X,-mK_X)\ne0\) for some integer \(m>0\), or
equivalently whether some positive multiple of \(-K_X\) is linearly
equivalent to an effective divisor.

We prove the following nonvanishing theorem.

\begin{theorem}[Anticanonical nonvanishing]\label{thm:headline}
Let \(X\) be a smooth connected projective variety over \(\C\).
If \(-K_X\) admits a smooth Hermitian metric of semipositive curvature,
then \(H^0(X,-mK_X)\ne0\) for some integer \(m>0\).
\end{theorem}

This is the positive-multiple formulation associated with Yau's
anticanonical nonvanishing question \cite[Problem~75, p.~395]{Yau1994}.
The multiple matters even when the anticanonical bundle is not torsion.
For a complex Enriques surface \(E\), \(K_E\) is nontrivial of order two
and has no section \cite[Section~3]{Dolgachev2016}.
Write \(p_E\) and \(p_{\mathbb P^1}\) for the two projections from
\(E\times\mathbb P^1\). Its anticanonical line is
\[
 p_E^*(-K_E)\otimes p_{\mathbb P^1}^*\OO_{\mathbb P^1}(2).
\]
The flat metric on the first factor and Fubini--Study metric on the
second make it smoothly semipositive. The product formula for sections
gives no first-power section, whereas the square has sections.
Restriction to a \(\mathbb P^1\) fiber shows that this anticanonical
line is not torsion.

\subsection{Geometry, monodromy, and earlier methods}

The structure theory explains both the strength of the metric assumption
and the difficulty that remains. Yau's prescribed-Ricci theorem
\cite{Yau1978} realizes the given smooth semipositive anticanonical
curvature as the Ricci form of a K\"ahler metric. The structure results
of Demailly--Peternell--Schneider \cite[Structure Theorem, p.~218]{DPS1996} and
Campana--Demailly--Peternell \cite[Theorem~1.4]{CDP2015} describe the
universal cover by its flat, compact Ricci-flat, and remaining compact
factors. The remaining factor is projective and has no holomorphic forms
of positive degree. After passage to a finite-index deck subgroup, the
actions that remain on this factor lie in a compact torus.

This last action cannot be ignored. A section on the compact factor
descends only if it is compatible with the deck transformations.
The flat and Ricci-flat factors supply canonical frames; on the remaining
factor we construct an anticanonical section fixed by the torus, for the
natural action on that factor's anticanonical bundle. Its product with those frames descends
to a finite \'etale cover, and its finite \'etale norm then
gives a section on \(X\). This reasoning uses the action on the universal
cover, not a global product decomposition of a finite cover.

In the broader nef setting,
Lazi\'c--Matsumura--Peternell--Tsakanikas--Xie proved numerical effectivity
for projective log canonical threefold pairs with rational boundary
and nef log-anticanonical divisor, under
\(\mathbb Q\)-factoriality or rational singularities
\cite[Theorem~A]{LMPTX2023}. Numerical effectivity means that the divisor
class is numerically equivalent to an effective rational divisor;
the existence of a section requires rational \emph{linear} effectivity.
For projective klt \(\mathbb Q\)-pairs with nef log-anticanonical
divisor, M\"uller proved actual nonvanishing in dimension three,
and in arbitrary dimension when the log-anticanonical divisor of a
general fiber of the maximal rationally connected quotient is semiample
\cite[Theorem~A and Corollary~B]{Muller2025}.
His equivariant theorem constructs naturally invariant plurisections
of a semiample log-anticanonical divisor for boundary-preserving actions
of commutative linear algebraic groups on projective sub-log-canonical
\(\mathbb Q\)-pairs
\cite[Theorem~C]{Muller2025}. It is a direct methodological predecessor
for addressing residual monodromy. Here the finite-volume argument
constructs the needed naturally invariant section from smooth
semipositivity, without assuming that fiber semiampleness.

In dimension two, Chen--Filip--Sun--Tosatti--Zhang have classified
compact K\"ahler surfaces with Hermitian-semipositive anticanonical bundle
\cite[Theorem~1.3]{CFSTZ2026}.
Their classification gives complementary low-dimensional geometry;
the present argument develops a section-producing mechanism in arbitrary
dimension.\footnote{The authors report that ChatGPT~5.6~Sol suggested
a preliminary version of their Lemma~5.2, which they reformulated and
developed. This records their declared assistance, not an attribution
of authorship.}

The ordinary part of the proof builds on two further developments.
Generic positivity and foliation methods control cotangent tensors:
the relevant inputs come from movable-cone duality
\cite{BDPP2013}, movable slope theory \cite{GKP2016}, and
Campana--P\u aun's positive-slope algebraicity criterion
\cite[Theorem~1.1]{CampanaPaun2019}.
The slope argument follows the foliation strategy used by
Ou \cite[Theorem~1.4]{Ou2023} and by
Lazi\'c--Matsumura--Peternell--Tsakanikas--Xie
\cite[Theorem~4.1]{LMPTX2023}.
Here the relative Bergman metric supplies ramification-corrected
positivity when an effective boundary and a finite adjoint measure
replace nef anticanonical data.
Hard Lefschetz with multiplier ideals \cite[Theorem~2.1.1]{DPS2001}
then produces twisted forms, whose determinants are the effective
divisors used to select a smaller base.

The transfer in Section~\ref{sec:transfer} uses the relative Bergman metric of
Berndtsson--P\u aun \cite[Theorem~0.1]{BerndtssonPaun2008}.
The ordinary proof uses its singular-coefficient form. The equivariant
proof instead uses Berndtsson's smooth direct-image positivity
\cite[Theorem~1.2]{Berndtsson2009}: averaging is an orthogonal holomorphic
projection onto the invariant line. These are distinct applications of
positivity, and we state their hypotheses separately.

\subsection{A finite-volume theorem and the proof strategy}
\label{sec:overview}

The inductive statement involves two line bundles. We use additive
notation: $J+L$ means $J\otimes L$, and $mL$ means $L^{\otimes m}$.
For an effective integral divisor $D$, write $s_D$ for the canonical
section of $\OO_Z(D)$. If $J=-K_Z+D$, this section is a $J$-valued
holomorphic top form. A metric $h_J$ therefore pairs $s_D$ with its
conjugate to give a measure, denoted $|s_D|_{h_J}^2$.

A singular metric has local squared frame norm $e^{-\phi}$. It is
semipositive when $\phi$ is plurisubharmonic; locally bounded weights
means bounds on both sides.

\begin{theorem}[Finite volume and sections with prescribed poles]
\label{thm:finite-volume}
Let $Z$ be a smooth connected projective complex variety of dimension
$d$ with $H^0(Z,\Omega_Z^q)=0$ for every $1\le q\le d$.
Let $D\ge0$ be an integral divisor. Suppose that $J=-K_Z+D$ has a
semipositive singular Hermitian metric $h_J$ satisfying
\[
                 \int_Z |s_D|_{h_J}^2<\infty.
\]
If a line bundle $L$ on $Z$ has a semipositive singular Hermitian metric
with locally bounded weights, then there exist integers $a>0$ and
$b\ge0$ such that
\[
                       H^0(Z,aL+bD)\ne0.
\]
\end{theorem}

This generality is needed because a fibration need not transfer the
line whose section is sought to the anticanonical bundle of its base.
Its fiber integrals instead produce a new target line on the base and
a separate corrected anticanonical line carrying the volume metric.
The volume metric $h_J$ may be unbounded, provided the indicated
measure is integrable. The metric on $L$ has the separate boundedness
assumption. The conclusion says that a positive power of $L$
has a rational section with poles only on $D$; the pole order $b$ is not
uniform. Rationally connected varieties satisfy the no-forms hypothesis,
but the theorem also applies more generally. For example, an Enriques
surface has no positive-degree holomorphic forms and a torsion canonical
bundle, although it is not rationally connected.

Theorem~\ref{thm:finite-volume} is proved by induction on dimension.
We first describe how its components fit together.
Assume that its conclusion fails. Then $L+D$ is pseudoeffective but not
big. A supporting movable curve class $\alpha\ne0$ satisfies
$L\cdot\alpha=D\cdot\alpha=0$. Finite adjoint volume bounds the slopes
of cotangent tensors against $\alpha$; in particular their determinant
lines have nonpositive degree. This is the content of
Section~\ref{sec:slope}.

The absence of holomorphic forms gives $\chi(\OO_Z)=1$. Riemann--Roch
and hard Lefschetz then produce sections
\[
 s_i\in H^0(Z,M+m_iL),\qquad m_i\longrightarrow\infty,
\]
where $M$ is one fixed cotangent determinant with $M\cdot\alpha\le0$.
Consider all nonempty complete systems $|uL+vM+wD|$ with nonnegative
integral coefficients. None can define a generically finite map: that
would make its divisor big, contradicting its nonpositive pairing with
$\alpha$. Choose a system with maximal image dimension and take the
relative algebraic closure of its image field in $\C(Z)$. This gives a
rational fibration onto a strictly smaller smooth projective base $S$.
Products of sections force every ratio in the indicated systems to be
a function on $S$.

This control of section ratios has two consequences. First, ratios among the
$s_i$ give a rational section $\rho$ of some $kL$, $k>0$, with no poles
on divisors dominating $S$. Second, on a smooth resolution $Y$ with maps $\pi:Y\to Z$ and
$f:Y\to S$, put $K_{Y/S}=K_Y-f^*K_S$. The relative adjoint spaces for
$K_{Y/S}+\pi^*(J+jkL)$ have dimension one for every $j\ge0$.
They are generated over the base field by the pulled-back form $s_D$
with its birational Jacobian,
divided by a base canonical frame and multiplied by $\rho^j$.
We prove transfer first under these explicit hypotheses;
Section~\ref{sec:induction} then constructs its inputs and closes the
induction, including the extension and common denominator arguments.

Section~\ref{sec:transfer} explains why these algebraic facts are the
right ones. On a normal equidimensional model, minima of vertical
divisor orders turn $\rho$ into a regular section $e$ after subtracting a
base line $A$. On a smooth resolution, integrate its powers against
the pulled-back adjoint measure. The zeroth integral gives a metric on
$-K_S+D_S$ with finite adjoint volume. Normalized high moments give a
locally bounded semipositive metric on $A$. Normalization leaves $e$
nonzero on some component above every base divisor. When integrating
$e^j$, the boundary estimate therefore has a fixed order, rather than
one growing with $j$; dividing its logarithm by $j$ removes that
boundary contribution. The correction $D_S$ is
supported on base primes for which every component upstairs is exceptional over
$Z$ or lies above $D$. Thus the induction hypothesis on $S$ yields a
section that lifts with no forbidden pole on $Z$.

To obtain the torus-invariant section needed by the geometric reduction,
Section~\ref{sec:torus} first proves the smooth invariant transfer
statement, Proposition~\ref{prop:invariant-transfer}, and then applies
it to a rational orbit space.
Fundamental vector fields give the relative section, and a dense orbit
gives invariant adjoint rank one in every degree. Smooth direct-image
positivity applies to the invariant line because compact averaging is
an orthogonal holomorphic projection. The ordinary theorem on the
quotient base then produces a naturally invariant anticanonical section.
The cover and norm results of Section~\ref{sec:structure} finish the
proof of Theorem~\ref{thm:headline}.

The proofs below establish the model preparation, slope bound, and
rank and transfer arguments using the cited geometric and analytic results.
Section~\ref{sec:models} introduces the two birational models and the
metric conventions used throughout.

\section{Finite covers and residual monodromy}
\label{sec:cover}\label{sec:structure}

A smooth semipositive anticanonical metric determines much of the
geometry of the universal cover.  The remaining issue for sections is
the action of the deck group.  We record the geometric reduction with
that action retained, so that an invariant section on the final compact
factor can later be descended to a finite cover of the original variety.
The geometric inputs are the Ricci-semipositive structure theorem of
Demailly--Peternell--Schneider \cite[Structure Theorem and \S3]{DPS1996}
and the rational-connectedness refinement of
Campana--Demailly--Peternell \cite[Theorem~1.4]{CDP2015}.

\begin{theorem}[Finite cover with compact torus monodromy]
\label{thm:structure}
Let $X$ be a smooth connected projective complex variety whose
anticanonical bundle has a smooth Hermitian metric of semipositive
Chern curvature.  There are a connected finite \'etale cover
$\nu:S\to X$, an integer $a\geq0$, and smooth connected projective
varieties $F$ and $Z$ with the following properties.
\begin{enumerate}
\item For a suitable K\"ahler metric on $X$, the common universal cover
of $X$ and $S$ has a holomorphic isometric decomposition
\begin{equation}\label{eq:cover-product}
             \widetilde X=\mathbb C^a\times F\times Z.
\end{equation}
The compact factors are simply connected.  The induced metric on $F$
is Ricci-flat, and $K_F$ has a nowhere-zero parallel holomorphic section.
The anticanonical bundle of $Z$ is smoothly semipositive, and
\begin{equation}\label{eq:cover-noforms}
 H^0(Z,\Omega_Z^p)=0\quad(p>0),\qquad
 H^q(Z,\mathcal O_Z)=0\quad(q>0).
\end{equation}
The variety $Z$ is rationally connected when it has positive dimension.

\item There are a full lattice $\Lambda\subset\mathbb C^a$ and a
homomorphism $\rho:\Lambda\to\operatorname{Aut}(Z)$ for which
\begin{equation}\label{eq:cover-action}
 S=(\mathbb C^a\times F\times Z)/\Lambda,\qquad
 \lambda\cdot(z,f,v)=(z+\lambda,f,\rho(\lambda)v).
\end{equation}
The closure $T$ of $\rho(\Lambda)$ in the holomorphic isometry group
of $Z$ is a compact connected real torus.  The trivial torus is allowed.

\item The action of $T$ extends to an algebraic action of a complex
algebraic torus $G$ on $Z$, with $T$ Zariski dense in $G$.
On $K_Z^{-1}=\det T_Z$ the action is the natural one induced by the
differential.  For every $m\geq0$, its representation on sections satisfies
\begin{equation}\label{eq:cover-invariants}
       H^0(Z,K_Z^{-m})^T=H^0(Z,K_Z^{-m})^G.
\end{equation}

\item The canonical bundles of $\mathbb C^a$ and $F$ have nowhere-zero
holomorphic frames $\eta_0$ and $\eta_F$ invariant under $\Lambda$.
For every integer $m\geq0$, the product projections
$\operatorname{pr}_0$, $\operatorname{pr}_F$, and $\operatorname{pr}_Z$
give the natural, $\Lambda$-equivariant identification
\begin{equation}\label{eq:cover-canonical}
 K_{\widetilde X}^{-m}
 =\operatorname{pr}_0^*K_{\mathbb C^a}^{-m}
   \otimes\operatorname{pr}_F^*K_F^{-m}
   \otimes\operatorname{pr}_Z^*K_Z^{-m},
\end{equation}
in which the first two factors have the invariant frame
$\eta_0^{-m}\otimes\eta_F^{-m}$.
\end{enumerate}
Any of the factors may be a point.  A point contributes the canonical
frame $1$; if $a=0$, the lattice is the zero group.
\end{theorem}

The quotient in \eqref{eq:cover-action} describes a locally trivial
bundle over $\mathbb C^a/\Lambda$.  Its monodromy $\rho$ can have infinite
image, so the finite cover $S$ need not be a product.  This is the reason
for retaining the natural torus action on the anticanonical bundle.

\begin{proof}
We first obtain the product and its cohomological properties.  We then
choose a finite-index deck subgroup and identify its compact action
with an algebraic torus action.

\smallskip\noindent
\emph{The metric splitting and the compact factors.}
The curvature form of the given metric on $K_X^{-1}$ is a smooth
semipositive representative of its first Chern class, with the usual
normalizing factor.  Yau's prescribed-Ricci theorem therefore gives a
K\"ahler metric on $X$ whose Ricci form is that curvature form
\cite{Yau1978}; this application is also stated explicitly in
\cite[p.~218]{DPS1996}.  In particular its Ricci tensor is nonnegative.
The structure theorem of \cite{DPS1996} gives a holomorphic isometric
decomposition
\[
          \widetilde X=\mathbb C^a\times\prod_i M_i,
\]
where each $M_i$ is compact and simply connected and the nonflat
factors are irreducible in the real de Rham decomposition.  The
compactness of these factors is the Cheeger--Gromoll input in the
structure theorem; the holomorphic decomposition comes from the
parallel complex structure.  Group the Ricci-flat factors into $F$
and the remaining factors into $Z$.

Here is the Bochner argument for the part of this decomposition used
below.  On a compact K\"ahler manifold with nonnegative Ricci tensor,
a holomorphic $p$-form $u$ is parallel.  Moreover, if
$\lambda_1,\ldots,\lambda_n\geq0$ are the Ricci eigenvalues in a
unitary frame and $u=\sum_{|J|=p}u_J\,dz_J$, the integrated Bochner
formula forces the pointwise vanishing of
\[
             \sum_{|J|=p}\left(\sum_{j\in J}\lambda_j\right)|u_J|^2.
\]
These are the Bochner consequences used in
\cite[\S2]{DPS1996}; equivalently see the formula in
\cite[(1.6)]{CDP2015}.  In particular, contraction of $u$ with a
positive Ricci eigendirection is zero.  If $u\ne0$ on an irreducible
factor, its contraction kernel
\[
             \{v\in T^{1,0}M_i:\iota_vu=0\}
\]
is a proper parallel complex subbundle.  At any point where the Ricci
tensor is nonzero, that kernel is also nonzero.  Its underlying real
subbundle contradicts irreducibility.  Thus every non-Ricci-flat
factor has no nonzero holomorphic form of positive degree.  The
decomposition of exterior forms on a product and K\"ahler Hodge
symmetry give \eqref{eq:cover-noforms}.  Rational connectedness of
these factors, including the irreducible compact Hermitian-symmetric
case, is the additional conclusion of
\cite[Theorem~1.4(a) and \S4]{CDP2015}; it is not inferred merely from
the vanishing of forms.

The canonical Chern connection of the Ricci-flat factor $F$ is flat.
Since $F$ is simply connected, parallel transport gives a nowhere-zero
parallel holomorphic frame $\eta_F$ of $K_F$.  The metric induced on
$K_Z^{-1}$ has semipositive curvature, since the Ricci tensor of the
product restricts to that of $Z$.

Projectivity of each compact factor follows directly from that of
$X$.  Fix points in the other factors and compose its slice in
$\widetilde X$ with the covering map to $X$.  This holomorphic map
has injective differential.  Pulling back a positive Hermitian metric
on an ample line bundle of $X$ therefore gives a positive line bundle
on the compact factor.  Kodaira's embedding theorem makes that factor
projective \cite{Kodaira1954}.  The argument requires the slice map
to be immersive, not injective.

\smallskip\noindent
\emph{Choosing the deck subgroup.}
Let $\Gamma=\pi_1(X)$ act on $\widetilde X$ by deck transformations.
These are holomorphic isometries.  Uniqueness of the de Rham
decomposition implies that, after a finite-index passage removing
permutations of its factors, every deck transformation is a product
map.  Put $M=F\times Z$.  Projection of this deck subgroup to the
Euclidean isometry group is discrete and has finite kernel.  Indeed,
for a fixed bound on the displacement of $0\in\mathbb C^a$, a
sufficiently large closed Euclidean ball $B$ satisfies
\[
       \gamma(B\times M)\cap(B\times M)\ne\varnothing
\]
for every deck transformation within that bound.  Since $B\times M$
is compact, proper discontinuity allows only finitely many such
transformations.  This proves both assertions.  The Euclidean image
is cocompact, by compactness of the quotient of the full product.
Bieberbach's theorem supplies a finite-index translation lattice in
this image, as in \cite[\S3(ii)]{DPS1996}.

For completeness, the finite kernel does not obstruct choosing an
actual lattice subgroup of the deck group.  Above the translation
lattice we have an extension
\[
       1\longrightarrow N\longrightarrow\Gamma_1
        \longrightarrow\Lambda_1\longrightarrow 0
\]
with $N$ finite.  Pass to the centralizer of $N$ in $\Gamma_1$;
this has finite index because $\operatorname{Aut}(N)$ is finite.
Its kernel $N'$ is central, and its image $\Lambda'$ is a sublattice
of finite index in $\Lambda_1$.  Choose lifts
$g_1,\ldots,g_{2a}$ of a lattice basis $\lambda_1,\ldots,\lambda_{2a}$ of $\Lambda'$.
Their commutators belong to $N'$.  If $e>0$ is divisible by the orders
of all elements of $N'$, then
\[
                  [g_i^e,g_j^e]=[g_i,g_j]^{e^2}=1.
\]
The commuting elements $g_i^e$ generate a free abelian subgroup:
a relation among them projects to a relation among the independent
vectors $e\lambda_i$.  Its image is $e\Lambda'$, so this subgroup
has finite index in the original deck group.  In case $a=0$, the deck
group is finite and we instead take its trivial subgroup.

We next eliminate the action on $F$.  The holomorphic isometry group
of a compact simply connected Ricci-flat manifold is finite.  To see
this, its Lie algebra consists of Killing fields.  The integrated
Bochner identity makes each such field parallel.  Its dual one-form
is closed and, by simple connectedness, exact.  An exact parallel
one-form on a compact manifold is zero: a primitive attains a maximum,
where its differential vanishes.  Thus the compact isometry group has
zero-dimensional Lie algebra and is finite.  Passing to the kernel of
the action on $F$ preserves a finite-index lattice subgroup.

The remaining action on $Z$ has commuting image in its compact
holomorphic isometry group.  Its closure is a compact abelian Lie
group.  Its identity component is a torus and its component group is
finite.  Pass once more to the preimage of that identity component.
The new image is dense in the identity component, since the latter is
open in the old closure.  Denote this connected closure by $T$, the
deck lattice by $\Lambda$, and its action on $Z$ by $\rho$.  This
proves \eqref{eq:cover-action}.

The corresponding quotient $S$ is a connected finite unramified
holomorphic cover of $X$.  The pulled-back positive line bundle
makes $S$ projective, and the covering is a finite \'etale
morphism of projective varieties.  We have now accounted for every
part of the deck action.  In particular $\eta_F$ is invariant because
the action on $F$ is trivial, while
$\eta_0=dz_1\wedge\cdots\wedge dz_a$ is invariant under translations.
The determinant of the product tangent decomposition gives
\eqref{eq:cover-canonical}, with its natural equivariance.

\smallskip\noindent
\emph{The algebraic torus.}
It remains to put the compact action on $Z$ in the algebraic category.
The point case is immediate, so suppose $\dim Z>0$ and choose a very
ample line bundle $A$ on $Z$.  By \eqref{eq:cover-noforms}, the
exponential sequence makes
\[
                 c_1:\operatorname{Pic}(Z)\longrightarrow
                      H^2(Z,\mathbb Z)
\]
injective.  Every element of the connected group $T$ is isotopic to
the identity and hence acts trivially on integral cohomology.
Consequently $t^*A\simeq A$ for every $t\in T$.
The complete linear system of $A$ therefore defines a canonical
projective representation
\[
           T\longrightarrow
           \operatorname{PGL}\bigl(H^0(Z,A)^*\bigr)
\]
preserving the embedded variety $Z$.  Choices of isomorphisms
$t^*A\simeq A$ change the associated linear maps only by scalars,
so this projective representation is well-defined.
It is continuous: choose a projective frame from the nondegenerate
irreducible embedded variety; a projective transformation is
determined continuously by the images of that frame.  Such a frame
exists by first choosing a projective basis and then a point outside
its finitely many coordinate hyperplanes.

Write $r=\dim H^0(Z,A)$.  To diagonalize the compact projective
action, take its compact preimage under the finite covering
$\operatorname{SL}_r(\mathbb C)\to\operatorname{PGL}_r(\mathbb C)$.
The identity component of this preimage is abelian and surjects onto
$T$: its Lie algebra maps isomorphically to the abelian Lie algebra
of $T$, and its image is an open and closed subgroup of $T$.
Averaging a Hermitian form makes this connected preimage unitary;
its commuting matrices can then be diagonalized simultaneously.
Thus, after conjugation, $T$ lies in a diagonal algebraic torus in
$\operatorname{PGL}_r(\mathbb C)$.

Let $G$ be its Zariski closure there.  It is connected, since a
connected $T$ cannot meet distinct components of an algebraic group
without being disconnected.  A connected algebraic subgroup of a
diagonal torus is an algebraic torus.  The projective stabilizer of
$Z$ is Zariski closed, so $G$ preserves $Z$ and acts algebraically.
Differentiation gives the natural algebraic linearization on
$K_Z^{-1}$.  The induced finite-dimensional representation on
$H^0(Z,K_Z^{-m})$ is algebraic.  The stabilizer of any section is
Zariski closed, and Zariski density of $T$ proves
\eqref{eq:cover-invariants}.  All the constructions respect point
factors and the stated canonical-frame conventions.
\end{proof}

\subsection{The norm of a section}

Once a section has been constructed on the finite cover, the following
elementary norm returns it to the original variety.  Its exponent is
the degree of the cover; no uniform degree is needed in the application.

\begin{lemma}[Finite \'etale norm]
\label{lem:etale-norm}
Let $\nu:S\to X$ be a finite \'etale morphism of degree $d>0$
between smooth connected projective complex varieties, and let $A$ be
a line bundle on $X$.  Every nonzero section
$s\in H^0(S,\nu^*A)$ has a nonzero norm
\[
                    \operatorname{Nm}_{\nu}(s)
                    \in H^0(X,A^{\otimes d}).
\]
In particular, a nonzero section of $K_S^{-m}$, for $m>0$, gives a
nonzero section of $K_X^{-md}$.
\end{lemma}

\begin{proof}
Take an analytic open set $U\subset X$ on which $A$ has a frame $e$
and the cover splits into $d$ sheets.  On these sheets write
$s=f_i\nu^*e$ for holomorphic functions $f_1,\ldots,f_d$.
Define the local norm by
\[
                    \left(\prod_{i=1}^d f_i\right)e^{\otimes d}.
\]
This expression is unchanged by a permutation of the sheets.
If the frame is replaced by $e'=u e$, the functions become
$f_i'=f_i/u$; their product changes by $u^{-d}$ and the new tensor
frame changes by $u^d$.  The expressions therefore glue to a
holomorphic section of $A^{\otimes d}$.

Since $S$ is smooth and connected it is integral.  The zero set of the
nonzero section $s$ is thus proper and has dimension less than
$\dim S$ (or is empty).  Its image under the finite map is a proper
closed subset of $X$.  Away from that image every sheetwise factor is
nonzero, so the norm is nonzero.  Projectivity identifies this
holomorphic section with an algebraic section.  Finally,
$K_S=\nu^*K_X$ because $\nu$ is \'etale; applying the result
with $A=K_X^{-m}$ proves the last assertion.  In dimension zero the
same construction is the product of nonzero constants.
\end{proof}

\section{Measures and models for rational fibrations}
\label{sec:models}

The induction will transfer two metrics along a rational map. Divisor
orders are best controlled on a normal equidimensional model; integration
requires a smooth model. We first record how to retain both models at
once. We also specify the analytic conventions behind the finite-volume
hypothesis.

\subsection{Adjoint measures and extension of weights}

In a holomorphic frame $v$ of a line bundle $J$, a semipositive singular
metric has squared frame norm $|v|^2_h=e^{-\phi}$ with $\phi$
plurisubharmonic (psh) and not identically $-\infty$. In particular
$\phi$ is locally bounded above and $|v|_h^2$ has a positive local
lower bound. A metric with locally bounded weights has bounds on both
sides. Pullbacks and products of the metrics used below preserve
semipositivity.

For $J=-K_Z+D$, write the $J$-valued top form $s_D$ in coordinates as
\[
 s_D=g(z)\,dz_1\wedge\cdots\wedge dz_d\otimes v.
\]
Its squared norm is the measure
\begin{equation}\label{eq:adjoint-measure}
       |s_D|_h^2=|g(z)|^2e^{-\phi(z)}d\lambda_z,
\end{equation}
with one fixed dimensional normalization of Lebesgue measure. The
expression is independent of coordinates and frame. Its total mass
need not be finite merely because the curvature is semipositive.
For a birational morphism $\pi:Y\to Z$ between smooth varieties,
the pulled-back form includes the Jacobian:
\[
 \pi^*s_D\in H^0(Y,K_Y+\pi^*J),\qquad
 \operatorname{div}(\pi^*s_D)=K_{Y/Z}+\pi^*D.
\]
Here $K_{Y/Z}=K_Y-\pi^*K_Z$ is the effective exceptional Jacobian
divisor. Change of variables preserves the total mass in
\eqref{eq:adjoint-measure}. All assertions about singular metrics and
fiber integrals will be made as assertions about measures or almost
everywhere defined functions, followed by their psh representatives.

We use two removable-singularity facts. A psh function on the
complement of an analytic set extends if it is locally bounded above.
Across an analytic set of codimension at least two in a smooth space,
this upper bound is automatic. To see the latter locally, project a
small polydisc in a generic complex direction. Boundaries of suitable
parallel discs lie in a compact set outside the analytic subset; almost
all such discs miss that subset. The maximum principle gives an upper
bound on those discs, and the submean inequality gives the same bound
at every point of the complement. Extension by upper limits is then
psh. These statements apply in line-bundle frames, and a uniform lower
bound survives extension by upper limits.

\subsection{A normal model and a smooth resolution}

\begin{lemma}[Prepared model]\label{lem:prepared-model}
Let $Z$ be smooth connected projective and $S_0$ an integral complex
variety. Let $Z\dashrightarrow S_0$ be a dominant rational map with connected
general fiber. There is a smooth projective birational model $S$ of
$S_0$ and a diagram
\[
\begin{tikzcd}[row sep=large,column sep=large]
Y\arrow[r,"r"]\arrow[dr,"f"']\arrow[rr,bend left=22,"\pi"]
 &\bar Y\arrow[r,"\bar\pi"]\arrow[d,"\bar f"]&Z\\
 &S&
\end{tikzcd}
\]
in which $\bar Y$ is normal projective, $Y$ is smooth projective,
$r$ and $\bar\pi$ are birational morphisms, and $\bar f$ is
surjective and equidimensional. The construction may start above a
prescribed birational model of $Z$. If an algebraic torus acts on $Z$
and the rational map is invariant, the diagram can be chosen
equivariant with trivial action on $S$.

Every prime divisor on $\bar Y$ either dominates $S$ or maps onto a
prime divisor of $S$. At a general point of the strict transform in
$Y$ of a prime $Q\subset\bar Y$ above a base prime $T$, there are
coordinates in which
\begin{equation}\label{eq:power-map}
 f(t,x_2,\ldots,x_d)=(t^e,x_2,\ldots,x_s),
 \qquad e=\ord_Q(\bar f^*T)>0,
\end{equation}
where $s=\dim S$ and $d=\dim Z$.
\end{lemma}

\begin{proof}
First replace $S_0$ birationally by a projective model of its function
field and compose the rational map with that replacement. This preserves
the generic fiber and any trivial base action.
Resolve the graph of the resulting rational map, starting above the prescribed
model if one has been given. Flatten the resulting projective morphism
by a modification of its base, using the strict transform
\cite{RaynaudGruson1971}. Resolve the modified base and take the main
component of the base change, namely the closure of the component over
the common dense open. Before taking that component, the flat family
has all fibers of the generic dimension $d-s$; base change preserves
this. A closed component has fiber dimension at most $d-s$.
Normalize the main component to obtain $\bar Y$. Normalization is
finite, so this upper bound persists. The dimension inequality for a
dominant map gives the reverse bound for every irreducible fiber
component. Thus $\bar f$ is equidimensional. Resolve $\bar Y$ to
obtain $Y$; no equidimensionality is asserted for $f$.

For an invariant rational map, use its invariant graph. The strict
transform and main component are invariant, the action lifts uniquely
through normalization, and characteristic-zero equivariant resolution
supplies the smooth source. All base modifications have trivial torus
action, so the resulting maps are equivariant.

If a prime $Q$ has image of codimension at least two in $S$, its generic
fiber over its image has dimension at least
$(d-1)-(s-2)=d-s+1$, contradicting equidimensionality. Hence the only
vertical divisors lie above base divisors. At general points of $Q$
and $T$, both divisors are smooth, $Y\to\bar Y$ is an isomorphism,
and the induced map $Q\to T$ is a submersion. A local equation of $T$
pulls back to $t^e$ times a unit. Take a holomorphic $e$th root of that
unit and absorb it into $t$; choose the remaining base coordinates
among coordinates on $Q$. This gives \eqref{eq:power-map}. Generic
smoothness and constructibility ensure that these neighborhoods occur
over a dense open of $T$, while avoiding any specified proper closed
subset of $Q$.
\end{proof}

A divisor dominating the base is called \emph{horizontal}; the others
are \emph{vertical}. The distinction in Lemma~\ref{lem:prepared-model}
will be used twice. Saturating pulled-back base differentials removes
the factor $t^{e-1}$ in \eqref{eq:power-map}; this is the ramification
correction in the slope proof. Later, the same coordinates estimate the
fiber integrals at base divisors. Exceptional divisors created on the
smooth resolution can map into higher codimension. They cause no loss
of divisor control because normalization is performed first on
$\bar Y$, and the final psh extensions are carried out on the smooth
base $S$.

\section{Cotangent tensors and the permitted boundary}
\label{sec:h:slope}\label{sec:slope}

The induction below produces effective divisors of the form $M+mL$,
where $M$ is a line bundle mapping nontrivially into a cotangent tensor
power.  We must
control the contribution of $M$ using the same boundary that is allowed
in the eventual section.  The finite adjoint volume gives precisely
this control.  The first step is analytic: a relative Bergman metric
remains positive after removing the vanishing of the differential along
multiple fibres.  The second step applies this observation to the
positive part of the tangent sheaf's slope filtration.

Throughout this section, $Z$ is a smooth connected projective complex
variety, $D\geq0$ is an integral divisor, and
\[
                        J=-K_Z+D.
\]
We assume that $J$ has a singular Hermitian metric $h_J$ of semipositive
curvature and that
\begin{equation}\label{eq:h:slope:volume}
              \int_Z |s_D|_{h_J}^2<\infty.
\end{equation}
Here $s_D$, considered as a section of $K_Z+J$, is a $J$-valued top
form, and the norm denotes its associated positive density.  A local
weight $\phi$ for $h_J$ satisfies $|e|_{h_J}^2=e^{-\phi}$ in a
holomorphic frame $e$.  It is plurisubharmonic, so the frame norm has a
strictly positive lower bound on each relatively compact chart.  No
upper bound for that norm is assumed.

\subsection{Removing the ramification from the Bergman metric}

A saturated foliation $\mathcal F\subset T_Z$ is a coherent subsheaf
with torsion-free quotient, closed under the Lie bracket.  Write
\[
 Q=T_Z/\mathcal F,\qquad
 P=\det Q^*=(\bigwedge^{\operatorname{rank}Q}Q^*)^{**}.
\]
Since $Z$ is smooth, $P$ is a line bundle.  If the foliation is induced
by a rational fibration, then $P$ agrees generically with the pullback
of the base canonical bundle.  They can differ along a multiple fibre;
that difference is the point of the following lemma.

\begin{lemma}\label{h:slope:foliation}
Under the assumptions above, let $\mathcal F\subset T_Z$ be an
algebraically integrable saturated foliation with
$0<\operatorname{rank}\mathcal F<\dim Z$.  Then the line bundle
\[
                    D-P=D+\det Q
\]
admits a singular Hermitian metric of semipositive curvature.  In
particular, $c_1(Q)+[D]$ is pseudoeffective.
\end{lemma}

\begin{proof}
Set $n=\dim Z$ and $b=\operatorname{rank}Q$.  Resolve the rational
fibration using Lemma~\ref{lem:prepared-model}.  Thus there are
morphisms
\[
 r:Y\longrightarrow\bar Y,\qquad
 \bar\pi:\bar Y\longrightarrow Z,\qquad
 \bar f:\bar Y\longrightarrow S,
 \qquad \pi=\bar\pi\circ r,\quad f=\bar f\circ r,
\]
where $Y,S$ are smooth projective, $\bar Y$ is normal, $\bar\pi$ and
$r$ are birational, and $\bar f$ is equidimensional with connected
general fibres.  Their dimension is $n-b$.  The foliation is the
saturated relative tangent sheaf at the generic point.

Pull back $s_D$ as a top form with coefficients, and denote the result
by
\[
               \eta\in H^0(Y,K_Y+\pi^*J).
\]
Its divisor is $K_{Y/Z}+\pi^*D$; the relative canonical divisor is the
effective birational Jacobian.  Change of variables off the exceptional
locus shows that the density of $\eta$ for $\pi^*h_J$ has finite total
mass.  This assertion concerns almost-everywhere densities, so no
pointwise product on the metric's polar set is needed.

Over a coordinate chart in the base, let $ds$ be a nonvanishing frame
of $K_S$.  On a smooth fibre $Y_s$, the restriction of
$\eta/f^*ds$ is a section of $K_{Y_s}+\pi^*J|_{Y_s}$.  It is nonzero
on general fibres, and Fubini's theorem gives finite squared norm on
almost every such fibre.  Berndtsson--P\u{a}un's relative Bergman
kernel theorem therefore supplies a semipositive metric on
\[
                       K_{Y/S}+\pi^*J
\]
\cite[Theorem~0.1]{BerndtssonPaun2008}.  Over the smooth-fibre locus, in a holomorphic frame its weight
$\psi$ is the logarithm of the supremum of the squared absolute
coefficients of fibre sections having squared norm at most one.
Only sections integrable for the restricted twisting metric enter
that supremum.

We now compare this metric with the line bundle in the statement.
Remove from $Z$ a closed subset of codimension at least two so that
$\pi$ is an isomorphism over the remaining open set and the sheaves
$\mathcal F,Q$ are locally free there.  At the generic point of every
divisor of this open set the conormal inclusion is
\[
                         f^*\Omega_S^1\longrightarrow Q^*.
\]
For a horizontal divisor it has full rank without a divisorial zero:
the restriction of $f$ to that divisor is generically submersive in
characteristic zero.  Every other divisor under consideration maps
onto a base divisor.  Indeed its strict transform on $\bar Y$ is a
prime divisor, and equidimensionality excludes a base image of
codimension at least two.

Fix such a vertical divisor $E$.  At its general points there are
coordinates on $Y$ and $S$ in which
\begin{equation}\label{eq:h:slope:power-map}
 s_1=z_1^e,\qquad s_i=z_i\quad(2\leq i\leq b),
 \qquad e=\operatorname{ord}_E(f^*\{s_1=0\}).
\end{equation}
To obtain these coordinates, take a point where the reduced map
$E\to\{s_1=0\}$ is submersive, write $f^*s_1=z_1^e u$ with $u$ a
unit, and absorb a local $e$-th root of $u$ into $z_1$.  Complete the
other pulled-back base coordinates to a coordinate system.  In this
chart the saturated conormal bundle is spanned by
$dz_1,\ldots,dz_b$.  The determinant of its displayed inclusion is
$e z_1^{e-1}$.  Consequently its vanishing order is exactly $e-1$.

This same integer occurs in the Bergman metric.  Choose a frame $v$
of $\pi^*J$ and use
\[
  \xi=(dz_1\wedge\cdots\wedge dz_n)\otimes(f^*ds)^{-1}\otimes v
\]
as the frame of $K_{Y/S}+\pi^*J$.  If $a\xi$ is a fibre section,
then, up to a constant of absolute value one from the order of wedge
factors, its coefficient as an actual fibre top form is
\[
                       \frac{a}{e z_1^{e-1}}.
\]
This distinction between relative and fibre coefficients is also the
one used in the definition of the relative Bergman kernel
\cite[Section~1]{BerndtssonPaun2008}.

Choose two concentric coordinate polydiscs whose closures lie in this
chart.  For each fixed $(z_1,\ldots,z_b)$ with $z_1\ne0$, the larger
polydisc in the free variables $z_{b+1},\ldots,z_n$ lies in the same
smooth nearby fibre.  Its radius can be fixed independently of the
base point.  The twisting frame norm has a uniform positive lower
bound there.  The holomorphic submean inequality on that polydisc,
for points in the smaller one, therefore gives
\[
 \|a\xi\|_{L^2(Y_s)}^2
 \ \geq\ c\,|z_1|^{-2(e-1)}
       \int_{\text{free polydisc}}|a|^2\,d\lambda
 \ \geq\ c'|a(z)|^2|z_1|^{-2(e-1)},
\]
where $c,c'>0$ do not depend on the fibre section or the nearby base
point.  For squared norm at most one, this implies
\begin{equation}\label{eq:h:slope:bergman-order}
                    \psi\leq(e-1)\log|z_1|^2+C.
\end{equation}
The estimate holds on the smooth-fibre locus; any exceptional
almost-everywhere qualifications disappear for the psh representative
by the submean inequality.

On the submersion locus where $\pi$ is an isomorphism we have
$K_{Y/S}+\pi^*J=D-f^*K_S$.  The determinant computation above says
that a frame of $D-P$ maps to $z_1^{e-1}$ times a nonvanishing multiple
of a frame of $D-f^*K_S$.  The induced metric on $D-P$ thus has weight
\[
                         \psi-(e-1)\log|z_1|^2
\]
up to a pluriharmonic frame term.  This weight is psh off $E$ and is
locally bounded above by \eqref{eq:h:slope:bergman-order}; hence it
extends psh across $E$.  At a horizontal divisor there is no
ramification correction, and the original Bergman metric already
extends.  The coordinate argument applies over a dense open subset
of each divisor.  The remaining excluded subset has codimension at
least two in the smooth variety $Z$, across which psh weights extend
uniquely.  The extensions respect changes of frame and therefore
define a semipositive metric on the global line bundle $D-P$.
\end{proof}

\subsection{The slope bound and the face used by induction}

A movable class is an element of the closed movable cone
$\operatorname{Mov}_1(Z)\subset N_1(Z)_{\mathbb R}$.  On a smooth
projective variety this cone is dual to the pseudoeffective divisor
cone \cite{BDPP2013}.  For a torsion-free coherent sheaf $E$ of
positive rank and a nonzero movable class $\alpha$, write
\[
 \mu_\alpha(E)=\frac{c_1(E)\cdot\alpha}{\operatorname{rank}E},
 \qquad \mu_{\min,\alpha}(E),\quad\mu_{\max,\alpha}(E)
\]
for its slope and extremal Harder--Narasimhan slopes.  Equivalently,
the minimum is the infimum of slopes of its nonzero torsion-free
quotients, and the maximum is the supremum of slopes of its nonzero
subsheaves.  We use the Harder--Narasimhan filtration and tensor slope
formula for movable classes, including real boundary classes
\cite[Corollary~2.27 and Theorem~4.2]{GKP2016}.

\begin{theorem}[Finite-volume cotangent bound]
\label{h:slope:cotangent}\label{thm:finite-volume-slope}
Let $Z$ be smooth connected projective over $\C$, let $D\ge0$ be
an integral divisor, and suppose that $J=-K_Z+D$ has a semipositive
singular Hermitian metric $h_J$ with $\int_Z|s_D|_{h_J}^2<\infty$.
If $\dim Z>0$,
then every nonzero real movable class $\alpha$ satisfies
\begin{equation}\label{eq:h:slope:min}
                  \mu_{\min,\alpha}(T_Z)\geq-D\cdot\alpha.
\end{equation}
For every integer $k\geq0$ and every line bundle $M$ with a nonzero
morphism $M\to(\Omega_Z^1)^{\otimes k}$, the line bundle
\begin{equation}\label{eq:h:slope:line}
                             -M+kD
\end{equation}
is pseudoeffective.  More generally, if $E$ is a torsion-free subsheaf
of rank $r>0$ in $(\Omega_Z^1)^{\otimes k}$, then
$-\det E+krD$ is pseudoeffective.
\end{theorem}

\begin{proof}
Assume first that $\dim Z>0$, fix $\alpha$, and set
$\delta=D\cdot\alpha\geq0$.  Semipositivity of $J$ gives
\begin{equation}\label{eq:h:slope:total}
                         c_1(T_Z)\cdot\alpha\geq-\delta.
\end{equation}
Let $\lambda_1>\cdots>\lambda_t$ be the slopes of the
Harder--Narasimhan quotients of $T_Z$, with positive ranks
$r_1,\ldots,r_t$.  Suppose, for a contradiction, that
$\lambda_t<-\delta$.  The total degree of the pieces with
nonpositive slopes then satisfies
\begin{equation}\label{eq:h:slope:negative-degree}
       \sum_{\lambda_i\leq0}r_i\lambda_i
                     \leq r_t\lambda_t<-\delta.
\end{equation}
By \eqref{eq:h:slope:total}, some piece has positive slope.
Consequently the positive part $\mathcal F$ of the filtration is
nonzero and proper.  It is saturated, and
\[
 \mu_{\min,\alpha}(\mathcal F)>0,
 \qquad \mu_{\max,\alpha}(T_Z/\mathcal F)\leq0.
\]

The bracket modulo $\mathcal F$ is an $\mathcal O_Z$-linear morphism
$\bigwedge^2\mathcal F\to T_Z/\mathcal F$.  Since the target is
torsion-free, it factors through
$A=(\bigwedge^2\mathcal F)/\mathrm{torsion}$.  When
$\operatorname{rank}\mathcal F>1$, the tensor slope formula, duality,
and the quotient inequality give
$\mu_{\min,\alpha}(A)\geq2\mu_{\min,\alpha}(\mathcal F)>0$.
The maximum slope of the target is nonpositive, so this morphism
vanishes.  If $\operatorname{rank}\mathcal F=1$, then $A=0$.
Thus $\mathcal F$ is a foliation.  Campana--P\u{a}un's algebraicity
criterion now shows that it is algebraically integrable
\cite[Theorem~1.1]{CampanaPaun2019}.

Lemma~\ref{h:slope:foliation} gives
$c_1(T_Z/\mathcal F)\cdot\alpha\geq-\delta$.  Its left side is
exactly the sum in \eqref{eq:h:slope:negative-degree}, a
contradiction.  This proves \eqref{eq:h:slope:min}.  In particular the
argument covers the cases in which the positive part would be zero
or the whole tangent sheaf: the assumed failure of the inequality
has already excluded both cases.

By duality and tensor-product compatibility,
\[
 \mu_{\max,\alpha}((\Omega_Z^1)^{\otimes k})
                 =-k\mu_{\min,\alpha}(T_Z)
                 \leq kD\cdot\alpha
\]
for $k>0$.  A nonzero map from $M$ is injective, because its target is
torsion-free.  Thus $(kD-M)\cdot\alpha\geq0$ for every movable
$\alpha$, and divisor--curve duality proves
\eqref{eq:h:slope:line}.  For the sheaf $E$, the same inequality
applied to its slope gives
$c_1(E)\cdot\alpha\leq krD\cdot\alpha$, which proves the determinant
assertion.  Determinants are taken reflexively; removing
codimension-two singular loci does not change their first Chern
classes.

If $k=0$, the target is $\mathcal O_Z$.  A nonzero map from $M$
provides a nonzero section of $-M$, and a positive-rank subsheaf of
$\mathcal O_Z$ has rank one with effective negative determinant.
These arguments also cover dimension zero.  In that dimension there
is no nonzero cotangent tensor for $k>0$, and no minimum slope of a
zero tangent sheaf is invoked.
\end{proof}

The coefficient of $D$ in this theorem is sharp even for curves.
Indeed, on a smooth projective curve of genus at least two, take a
nonzero holomorphic one-form $\omega$ and put $D=\operatorname{div}(\omega)$.
Then $J=-K_Z+D$ is trivial and its flat metric satisfies
\eqref{eq:h:slope:volume}.  Taking $M=kK_Z$ gives equality in the
cotangent bound; replacing $kD$ by $cD$ with $c<k$ gives negative
degree and therefore cannot give a pseudoeffective class.

\begin{corollary}[A supporting class for the induction]
\label{h:slope:face}
Under the hypotheses of Theorem~\ref{thm:finite-volume-slope},
assume $\dim Z>0$ and let $L$ be a pseudoeffective line bundle for
which $L+D$ is not big.  There exists a nonzero real movable class
$\alpha$ with
\[
                       L\cdot\alpha=D\cdot\alpha=0.
\]
For this class $\mu_{\min,\alpha}(T_Z)\geq0$.  Suppose moreover that
$M\to(\Omega_Z^1)^{\otimes k}$ is nonzero and that effective divisors
$N_i\sim M+m_iL$ are given.  Then
\[
                  M\cdot\alpha=N_i\cdot\alpha=0
                         \quad\text{for every }i.
\]
Every finite nonnegative linear combination of $D,N_1,N_2,\ldots$
has zero intersection with $\alpha$ and is not big.
\end{corollary}

\begin{proof}
The class of $L+D$ is on the boundary of the closed pseudoeffective
cone, whose interior is the big cone.  A supporting hyperplane and
cone duality give a nonzero real movable class $\alpha$ annihilating
it.  Both $L$ and $D$ pair nonnegatively with $\alpha$, so they pair
with it individually to zero.  Theorem~\ref{h:slope:cotangent} gives
the tangent inequality and $M\cdot\alpha\leq0$.  Effectivity of each
$N_i$ gives the reverse inequality through
$N_i\cdot\alpha=M\cdot\alpha$.  Finally an interior point of a
full-dimensional closed convex cone pairs strictly positively with
every nonzero member of its dual cone; hence no indicated combination
is big.
\end{proof}

The supporting class in this corollary can lie on the boundary of the
movable cone and need not be rational or represented by a single
moving curve.  Its role is numerical: it rules out a generically
finite linear system made from the displayed effective divisors and
thereby forces the dimension drop in the subsequent argument.

\section{Transferring a finite measure and a bounded metric}
\label{sec:transfer}

The next proposition isolates the induction step. It does not assume
the existence of a section on the base: it constructs base data to
which a lower-dimensional nonvanishing theorem can be applied.
We retain two models of the total space because they do different jobs.
On a normal equidimensional model every vertical prime maps onto a
base prime, so divisor orders can be normalized there. A smooth
resolution carries canonical forms and fiber integrals, but need not
remain equidimensional. Such diagrams arise from flattening a resolved
rational map \cite{RaynaudGruson1971}.

A prime is called \emph{horizontal} if it dominates the base, and
\emph{vertical} otherwise. For a morphism \(f:Y\to S\) of smooth varieties
we write \(K_{Y/S}=K_Y-f^*K_S\). Its restriction to a smooth fiber is the
canonical line of that fiber.

\begin{proposition}[Two-metric transfer]\label{prop:transfer}
Let \(X,S,Y\) be smooth connected projective varieties over \(\C\), and
let \(\bar Y\) be a normal projective variety. Suppose
\(r:Y\to\bar Y\) and \(\bar\pi:\bar Y\to X\) are projective birational
morphisms, and \(\bar f:\bar Y\to S\) is a surjective equidimensional
morphism with connected general fiber. Put
\(\pi=\bar\pi r\) and \(f=\bar f r\), as in Figure~\ref{fig:models}.

Let \(D\) be an effective integral divisor on \(X\), and let
\(J=-K_X+D\) have a semipositive singular metric \(h_J\) with
\[
 \int_X |s_D|_{h_J}^2<\infty.
\]
Let \(L\) be a line bundle on \(X\) with a semipositive metric \(h_L\)
of locally bounded weights. Suppose that for an integer \(k>0\)
there is a nonzero rational section \(\rho\) of \(k\bar\pi^*L\)
with no horizontal poles, and that
\begin{equation}\label{eq:rank}
 \dim_{\C(S)}
 H^0\!\left(Y_\eta,
  \left(K_{Y/S}+\pi^*(J+jkL)\right)|_{Y_\eta}\right)=1
 \qquad(j=0,1,2,\ldots).
\end{equation}
Here \(\eta\) is the generic point of \(S\) and
\(Y_\eta=Y\times_S\operatorname{Spec}\C(S)\).

After replacing \(k,\rho\) by \(dk,\rho^d\) for some integer \(d>0\),
there are a line bundle \(A\) on \(S\), an effective integral divisor
\(D_S\), and a nonzero regular section
\[
 e\in H^0(\bar Y,k\bar\pi^*L-\bar f^*A)
\]
with these properties:
\begin{enumerate}[label=\textup{(\roman*)}]
\item \(A\) has a semipositive metric with locally bounded weights.
\item \(-K_S+D_S\) has a semipositive singular metric \(h_S\) with
      \(\int_S|s_{D_S}|_{h_S}^2<\infty\).
\item \(D_S\) is supported on base primes \(T\) for which every prime
      component of \(\bar f^*T\) is \(\bar\pi\)-exceptional or lies in
      \(\Supp(\bar\pi^*D)\).
\item If \(a>0,b\ge0\) are integers and \(H^0(S,aA+bD_S)\ne0\), then
      \(H^0(X,akL+cD)\ne0\) for some integer \(c\ge0\).
\end{enumerate}
In these conclusions \(k\) denotes the enlarged integer.
\end{proposition}

\begin{figure}[ht]
\centering
\begin{tikzcd}[row sep=large,column sep=large]
Y \arrow[r,"r"] \arrow[dr,"f"'] \arrow[rr,bend left=22,"\pi"]
  & \bar Y \arrow[r,"\bar\pi"] \arrow[d,"\bar f"]
  & X\\
  & S &
\end{tikzcd}
\caption{The model diagram in Proposition~\ref{prop:transfer}.
The normal space \(\bar Y\) is equidimensional over \(S\); the smooth
space \(Y\) is used for integration. Both upper maps are birational.
The diagram does not assert that \(f\) is equidimensional.}
\label{fig:models}
\end{figure}

\begin{proof}
We first normalize \(\rho\) so that a divided section is regular and
nonzero on at least one component over each base prime. We then use
its adjoint integrals twice: the zeroth integral gives the volume
metric, and the high moments give the bounded metric. The support
of the first correction will make the final lifting possible.

If \(S\) is a point, every prime of \(\bar Y\) is horizontal.
Thus \(\rho\) is regular by normality and descends as a nonzero
section of \(kL\) on the smooth variety \(X\).
Take \(A\) trivial and \(D_S=0\); the metric assertions on a point
and the lifting assertion follow. We assume \(\dim S>0\) below.

\medskip\noindent
\emph{Divisorial normalization.}
For a prime divisor \(T\subset S\), define
\begin{equation}\label{eq:minima}
 c_T=\min_{\bar f(Q)=T}
       \frac{\ord_Q(\rho)}{\ord_Q(\bar f^*T)},
\end{equation}
where \(Q\) ranges over the prime divisors of \(\bar Y\) mapping onto
\(T\). These are finitely many components for each \(T\), and the
denominators are positive integers. Only finitely many \(c_T\) are
nonzero. Replace \(\rho\) by a power and \(k\) by the corresponding
multiple to make every \(c_T\) integral, and set
\(A=\OO_S(\sum_Tc_TT)\). The rank condition \eqref{eq:rank} persists:
the new adjoint twists form a subsequence of the old ones and still
include \(j=0\).

Divide \(\rho\) by the pullback of the tautological rational section
of \(A\), obtaining a rational section \(e\) of
\(k\bar\pi^*L-\bar f^*A\).
Its order is nonnegative at horizontal primes by hypothesis and at
vertical primes by \eqref{eq:minima}. Equidimensionality ensures that
every vertical prime maps onto a base prime. Thus these are all prime
divisors of \(\bar Y\), and normality makes \(e\) regular.
For each base prime \(T\), a component attaining the minimum in
\eqref{eq:minima} has order zero in \(e\).
We also write \(e\) for its pullback to \(Y\).

\medskip\noindent
\emph{The fiber integrals and their curvature.}
Pullback of the \(J\)-valued top form includes the Jacobian:
\[
 \eta_Y=\pi^*s_D\in H^0(Y,K_Y+\pi^*J),\qquad
 \mu_Y=|\eta_Y|_{\pi^*h_J}^2.
\]
Change of variables gives finite total mass for \(\mu_Y\).
In particular this assertion concerns measures and does not require
a pointwise value for a singular metric on its polar set.

On a coordinate chart \(U\subset S\), take a canonical frame \(ds\)
of \(K_S\) and a frame \(v\) of \(A\). The section
\(e_v=e f^*v\) belongs to \(k\pi^*L\) over \(f^{-1}(U)\). For
\(j=0,1,2,\ldots\), set
\begin{equation}\label{eq:moments}
 u_j=(\eta_Y/f^*ds)e_v^j,\qquad
 V_j(s)=\int_{Y_s}|u_j|_{\pi^*(h_Jh_L^{jk})}^2,
 \qquad \Phi_j=-\log V_j .
\end{equation}
The section \(u_j\) belongs to \(K_{Y/S}+\pi^*(J+jkL)\).
On a smooth fiber it is a top form with coefficients in the indicated
twisting line; \eqref{eq:moments} pairs that form with its conjugate
to give a positive density. Properness, regularity of \(e_v\), and
boundedness of \(h_L\)'s weights bound \(|e_v|^2\) on inverse images
of relatively compact base charts. Since \(V_0\) is the density of
the finite measure \(f_*\mu_Y\), all \(V_j\) are finite and positive
almost everywhere on a suitable smooth-fibration open.

For clarity, changes of frames have the following effect:
\begin{equation}\label{eq:frames}
 v\mapsto gv,\quad ds\mapsto h\,ds
 \quad\Longrightarrow\quad
 \Phi_j\mapsto \Phi_j-j\log|g|^2+\log|h|^2 .
\end{equation}
This will give weights on two different lines: \(-K_S\) at \(j=0\),
and \(A\) after division by \(j\) and passage to the limit.

The relative Bergman theorem
\cite[Theorem~0.1]{BerndtssonPaun2008}
applies to the semipositive twisting metric
\(\pi^*(h_Jh_L^{jk})\): \eqref{eq:moments} supplies a nonzero
square-integrable adjoint section on almost every smooth fiber.
The theorem produces a semipositive relative Bergman metric.
On a dense base open where base change holds and \eqref{eq:rank}
has dimension one, \(u_j\) generates the adjoint section space.
For almost every such fiber \(V_j<\infty\), so its one-dimensional
holomorphic adjoint space is also its full space of square-integrable
sections: every section is a scalar multiple of \(u_j\).
Consequently the Bergman weight, in a holomorphic frame upstairs, is
\[
 \Phi_j\circ f+\log|\widehat u_j|^2
\]
almost everywhere, where \(\widehat u_j\) is the coefficient of \(u_j\).
Off the zero locus of that coefficient, its logarithm is
pluriharmonic, so subtracting it leaves a psh function upstairs.
In a local product chart, Fubini's theorem gives almost all constant
slices on which the displayed identity holds almost everywhere.
Restriction to any of these slices gives a psh representative of
\(\Phi_j\) on the base. For fixed base frames, these representatives
agree almost everywhere on overlaps, hence everywhere, and therefore patch.

The base-change open can depend on \(j\), so we now put all these
weights on one fixed open \(S^\circ\). Choose \(S^\circ\) such that
\(f\) is smooth there and every fiber has a point where both
\(\eta_Y\) and \(e\) are nonzero. A psh weight is locally bounded above;
therefore its squared frame norm has a positive local lower bound,
even when the metric is singular. On a small product chart around
such a point, this lower bound for \(h_J\), together with the
bounds for \(h_L\), gives
\begin{equation}\label{eq:common-open}
 V_j(s)\ge C_1C_2^j,\qquad \Phi_j(s)\le C+jC',
\end{equation}
locally on \(S^\circ\), with positive \(C_1,C_2\) independent of \(j\).
Removable singularities extend each psh representative across its
additional exceptional base loci inside \(S^\circ\).
The union of the exceptional measure-zero sets for the countably many
\(j\) still has measure zero; hence all moment comparisons can be used
simultaneously. Almost-everywhere upper bounds also hold for the psh
representatives by the submean inequality.

We have now obtained all the moment weights on a fixed open set.
To continue the induction, both kinds of weight must extend to the
boundary. Their estimates use different components of a singular fiber.

\medskip\noindent
\emph{The zeroth moment: finite volume with supported correction.}
Call a base prime \(T\) \emph{bad} if every component of
\(\bar f^*T\) is \(\bar\pi\)-exceptional or lies in
\(\Supp(\bar\pi^*D)\). Only finitely many primes are bad: each is the
image of a member of the finite list of exceptional and boundary
divisors on \(\bar Y\).

Take a prime \(T\) in the complement of \(S^\circ\), and a component
\(Q\) above it on \(\bar Y\). At general points of its strict transform
in \(Y\), there are coordinates
\[
 s_1=t^\ell,\qquad s_i=x_i\quad(2\le i\le\dim S),
 \qquad \ell=\ord_Q(\bar f^*T)>0 .
\]
Indeed, the map on the divisors is generically smooth, and a root of
a nonvanishing unit absorbs the unit in the first coordinate.
Write the coefficient of \(\eta_Y\) as \(t^q\) times a unit at such
a point, with \(q\ge0\). Division by \(ds\) contributes
\(t^{-(\ell-1)}\), so the relative top-form coefficient has magnitude
comparable to \(|t|^{q-\ell+1}\). The metric lower bound and integration
over a fixed patch in the free fiber coordinates give
\begin{equation}\label{eq:zero-bound}
 V_0(s)\ge C|s_1|^{2(q-\ell+1)/\ell}
\end{equation}
almost everywhere near a general point of \(T\).

If \(T\) is not bad, choose \(Q\) nonexceptional over \(X\) and
outside \(\Supp(\bar\pi^*D)\). Then the birational Jacobian and \(s_D\)
are generically nonzero on \(Q\), so \(q=0\). Thus \(-\log V_0\)
is locally bounded above there. At the bad primes choose sufficiently
large nonnegative integral coefficients for \(D_S\). Equation
\eqref{eq:zero-bound} makes the weights
\begin{equation}\label{eq:base-volume-weight}
 \Phi_0+\log|\widehat s_{D_S}|^2
\end{equation}
locally bounded above at general points of every boundary prime.
Here the hat denotes the coefficient in a local frame of \(\OO_S(D_S)\).
These weights are psh off the boundary, extend across its general
divisorial points, and then extend across the remaining analytic set
of codimension at least two. The frame rule \eqref{eq:frames} gives a
metric on \(-K_S+D_S\).

Pairing \eqref{eq:base-volume-weight} with \(s_{D_S}\) cancels the
added factor. The resulting adjoint measure equals \(f_*\mu_Y\)
almost everywhere, so has finite total mass. We have obtained the
volume metric and the required support of its correction.

\medskip\noindent
\emph{High moments: bounded weights on \(A\).}
For \(l=1,2,\ldots\), put on \(S^\circ\)
\begin{equation}\label{eq:envelope}
 \Psi=\lim_{l\to\infty}
      \left(\sup_{j\ge l}\frac{\Phi_j}{j}\right)^* ,
\end{equation}
where the star denotes upper-semicontinuous regularization in fixed
local frames. Equation \eqref{eq:common-open} gives uniform local
upper bounds. The regularized suprema are therefore psh and decrease.

Their limit does not collapse to \(-\infty\).
On any relatively compact base chart \(U\), including charts meeting
the boundary, properness and bounded weights of \(h_L\) give a single
constant \(C_U>0\) with
\begin{equation}\label{eq:uniform-lower}
 |e_v|_{\pi^*h_L^k}^2\le C_U,\qquad V_j\le C_U^jV_0 .
\end{equation}
The first bound holds on the whole inverse image of \(U\):
take a finite cover of its compact closure by line-bundle frames.
For almost every \(s\in U\cap S^\circ\), \(V_0(s)\) is finite and
positive. Dividing \eqref{eq:uniform-lower} by \(j\) after taking
negative logarithms gives
\(\Psi\ge-\log C_U\).
Thus \(\Psi\) is psh. Its frame transformation is exact, although
that of each normalized moment has a vanishing error. On a relatively
compact overlap, \eqref{eq:frames} changes \(\Phi_j/j\) by
\(-\log|g|^2+(\log|h|^2)/j\). The last term has absolute value
at most \(C/l\) for every \(j\ge l\). The two tail suprema,
and then their upper-semicontinuous regularizations, therefore differ
from one another plus \(-\log|g|^2\) by at most \(C/l\).
Taking the decreasing limit gives precisely
\(\Psi\mapsto\Psi-\log|g|^2\), the weight rule for \(A\).

To extend \(\Psi\), now choose a component \(Q\) attaining the
minimum in \eqref{eq:minima}. This need not be the component used
for the zeroth moment. Its defining property is that \(e\) is nonzero
generically on \(Q\). The local norm of \(e_v\) is therefore bounded
below on a suitable patch. The same power-map calculation as in
\eqref{eq:zero-bound}, with the factor \(e_v^j\), gives
\begin{equation}\label{eq:high-bound}
 V_j(s)\ge C_1C_2^j|s_1|^{2b_T}
\end{equation}
for a nonnegative integer \(b_T\) independent of \(j\).
For example, increasing a nonnegative integer above
\((q-\ell+1)/\ell\) suffices after shrinking the coordinate chart.
The absence of zeros of \(e\) on the chosen component is what prevents
an additional exponent proportional to \(j\).

Shrink the chart so that \(|s_1|<1\), and decrease \(C_1\) if needed
so that \(C_1\le1\). Put
\[
 B(s)=-\log C_1-2b_T\log|s_1|\ge0.
\]
For \(j\ge l\), the negative logarithm of \eqref{eq:high-bound} gives
\[
 \frac{\Phi_j(s)}{j}\le-\log C_2+\frac{B(s)}{j}
                  \le-\log C_2+\frac{B(s)}{l}.
\]
The right side is continuous off \(T\), so it also bounds the
upper-semicontinuous regularization of the tail supremum there.
Letting \(l\to\infty\) in \eqref{eq:envelope} gives
\(\Psi\le-\log C_2\) throughout the punctured chart.
Hence \(\Psi\) is locally bounded above near general points of
each boundary prime, and extends psh across them and then across
codimension two. The lower bound in \eqref{eq:uniform-lower} was
uniform on the whole base chart, so the extension retains it.
The extended psh weight is locally bounded above as well.
Thus the metric on \(A\) has locally bounded weights on both sides.

\medskip\noindent
\emph{Lifting the section.}
We have the two metrics needed on the base. It remains to verify
that the divisor used for finite volume creates only permitted poles
on \(X\). Regard a nonzero section of \(aA+bD_S\) as a rational section
of \(aA\) with poles only on \(D_S\). Pull it to the normal
equidimensional model \(\bar Y\) and multiply by \(e^a\).
The resulting nonzero rational section of \(ak\bar\pi^*L\) has poles
only above bad primes. By their definition every such prime divisor
of \(\bar Y\) is exceptional over \(X\) or lies above \(D\).

At every prime of \(X\) outside \(D\), the strict transform in
\(\bar Y\) therefore has no pole. The rational section downstairs
extends there; exceptional divisors give no additional obstruction
on the smooth, hence normal, variety \(X\).
Its remaining finitely many pole orders are cleared by multiplication
by \(s_D^c\) for a sufficiently large integer \(c\ge0\).
This gives a nonzero section of \(akL+cD\), as required.
\end{proof}

\section{Choosing the smaller base and completing the induction}
\label{sec:induction}

We now prove Theorem~\ref{thm:finite-volume}. The transfer proposition
requires a rational section without horizontal poles and adjoint rank
one in every nonnegative degree. The task of this section is to obtain
both conditions from failure of nonvanishing. A maximal linear system
will give a strictly smaller base; the same maximality will control all
ratios on its generic fiber.

\subsection{One fixed determinant with unbounded twists}

We first explain the passage from cohomology to a single line bundle.
The determinant construction is a standard way to extract effective
divisors from twisted differentials; compare the determinant method of
Lazi\'c--Peternell \cite[Lemma~4.1]{LP2018}. We include the
argument because its uniformity in arbitrarily far tails is needed
here, including when the numerical class of $L$ is zero.

\begin{lemma}\label{lem:determinant}
Let $Z$ be smooth connected projective of dimension $d>0$, with
$H^0(Z,\Omega_Z^q)=0$ for $1\le q\le d$. Let $L$ have a
semipositive metric with locally bounded weights, and let $\alpha$ be
a nonzero movable class satisfying $\mu_{\min,\alpha}(T_Z)\ge0$.
There are a line bundle $M$ with $M\cdot\alpha\le0$, strictly
increasing positive integers $m_i\to\infty$, and nonzero sections
\[
                     s_i\in H^0(Z,M+m_iL).
\]
The line $M$ is the determinant of a saturated subsheaf of
$\Omega_Z^p$ for some $0\le p\le d$.
\end{lemma}

\begin{proof}
Hodge symmetry gives $H^q(Z,\OO_Z)=0$ for $q>0$, and hence
$\chi(\OO_Z)=1$. Riemann--Roch makes $\chi(K_Z+mL)$ a polynomial
in the integer $m$, whose constant term is
\[
                  \chi(K_Z)=(-1)^d\chi(\OO_Z)=(-1)^d
\]
by Serre duality. It is therefore nonzero for all but finitely many
positive integers. Some fixed $q$ has $H^q(Z,K_Z+mL)\ne0$ for
unbounded positive $m$.

For each such $m$, the metric $h_L^m$ has trivial multiplier ideal:
its weight is locally bounded, so every holomorphic germ is locally
square integrable. Hard Lefschetz with multiplier ideals
\cite[Theorem~2.1.1]{DPS2001} gives a surjection
\[
 H^0(Z,\Omega_Z^{d-q}\otimes mL)
       \longrightarrow H^q(Z,K_Z+mL)
\]
by wedging with the $q$th power of a K\"ahler form. Choose nonzero
maps $-mL\to U$, where $U=\Omega_Z^{d-q}$, for an increasing
sequence of these exponents.

At the generic point of $Z$, take the subspace of $U_{\C(Z)}$
spanned by all chosen images after a given cutoff. These subspaces
form a decreasing chain in a finite-dimensional vector space. Their
dimensions eventually stabilize, so the subspaces themselves stabilize
to a nonzero subspace $W$. Let $\mathcal G\subset U$ be its saturated
extension. It is coherent: saturate the images of finitely many maps
spanning $W$. Every map in a sufficiently far tail factors through
$\mathcal G$, because $U/\mathcal G$ is torsion-free and the
composite vanishes generically.

Put $r=\operatorname{rank}\mathcal G$ and
$M=(\bigwedge^r\mathcal G)^{**}$. This is a line bundle on smooth
$Z$. For every cutoff, choose $r$ generically independent maps from
the tail. Their exterior product is a nonzero section of $M+mL$,
where $m$ is the sum of their exponents. As the cutoff tends to
infinity, so does $m$. The maps and their determinant extend across
codimension two into the line $M$. Selecting an increasing subsequence
gives the sections in the statement.

For $d-q>0$, antisymmetrization embeds $U$ in
$(\Omega_Z^1)^{\otimes(d-q)}$. The tensor slope formula and the
assumed tangent inequality give $\mu_{\max,\alpha}(U)\le0$,
and therefore $M\cdot\alpha\le0$. If $d-q=0$, then
$U=\OO_Z$ and its nonzero saturated subsheaf is $\OO_Z$ itself;
thus $M=\OO_Z$ and the same conclusion holds.
\end{proof}

\subsection{A maximal system controls every section ratio}

The next lemma is algebraic. It isolates the dimension drop and the
rank calculation from the analytic transfer. The nonpositive degrees
are used only to prevent the base from having the full dimension.

\begin{lemma}[Maximal base and adjoint rank]\label{lem:maximal-base}
Let $Z$ be smooth connected projective of positive dimension $d$.
Let $L,M$ be line bundles and $D\ge0$ an integral divisor. Suppose
there is a nonzero movable class $\alpha$ with
\[
 L\cdot\alpha=D\cdot\alpha=0,\qquad M\cdot\alpha\le0,
\]
and nonzero sections $s_i\in H^0(Z,M+m_iL)$ for strictly increasing
positive integers $m_i\to\infty$.
There is a smooth projective base $S$ with $\dim S<d$, together with
the prepared diagram of Lemma~\ref{lem:prepared-model} and an integer
$k>0$, for which:
\begin{enumerate}[label=\textup{(\roman*)}]
\item a nonzero rational section $\rho$ of $k\bar\pi^*L$ has no
      horizontal poles on $\bar Y$;
\item putting $J=-K_Z+D$, one has for every integer $j\ge0$
\begin{equation}\label{eq:induction-ranks}
 \dim_{\C(S)}H^0\!\left(Y_\eta,
   \left(K_{Y/S}+\pi^*(J+jkL)\right)|_{Y_\eta}\right)=1.
\end{equation}
\end{enumerate}
If $S$ is a point, assertion \textup{(i)} already gives a nonzero
section of $kL$ on $Z$.
\end{lemma}

\begin{proof}
Consider the nonempty complete linear systems of
\begin{equation}\label{eq:systems}
              B=uL+vM+wD,
       \qquad u,v,w\in\mathbb Z_{\ge0}.
\end{equation}
Their degrees against $\alpha$ are nonpositive. A divisor whose
complete system defines a generically finite map is big. But a big
class is an interior point of the pseudoeffective cone, and so pairs
strictly positively with every nonzero member of the dual cone.
Thus every system in \eqref{eq:systems} has image dimension less than
$d$.

Choose $B_0$ in this collection with maximal image dimension $s$, and
let $V$ be the closed image of its rational map. This maximum exists
because the possible dimensions are integers between $0$ and $d-1$.
Let $K$ be the relative algebraic closure of $\C(V)$ in $\C(Z)$;
it is a finite extension of $\C(V)$, since $\C(Z)$ is a finitely
generated field extension. Take a smooth projective model $S$ of $K$,
with a morphism to $V$ after resolution. Its dimension is $s<d$, and
the rational map $Z\dashrightarrow S$ has connected general fiber.

Every ratio of two nonzero sections in any one system
\eqref{eq:systems} belongs to $K$. Indeed, combine that system with
$|B_0|$ by taking products of sections. The resulting linear subsystem
of $|B+B_0|$ gives the joint rational map. If its section ratio were
transcendental over $\C(V)$, the joint image would have dimension
at least $s+1$, contrary to the choice of $B_0$. The ratio is therefore
algebraic over $\C(V)$ and lies in its relative algebraic closure.
We will use this \emph{section-ratio property} both for interpolation
and for rank one.

Resolve $|B_0|$ before applying Lemma~\ref{lem:prepared-model}.
The resulting smooth model carries an effective divisor $R_0$ and a
line bundle $H_0$ on $S$ such that
\begin{equation}\label{eq:fixed-moving}
                  \pi^*B_0=f^*H_0+R_0.
\end{equation}
Here $H_0$ is the pullback of $\OO_V(1)$. If $s>0$ it is big,
since $S\to V$ is generically finite. All additional modifications
preserve the displayed effective decomposition by pullback.

Set $k=m_2-m_1>0$ and take
$\rho=\bar\pi^*(s_2/s_1)$. For $i>2$, the numerator and denominator
in
\begin{equation}\label{eq:interpolation}
                  \frac{s_i^k s_1^{m_i-m_2}}
                       {s_2^{m_i-m_1}}
\end{equation}
are sections of the same line bundle of the form
\eqref{eq:systems}: both have $M$ coefficient $m_i-m_1$ and $L$
coefficient $m_2(m_i-m_1)$. Thus \eqref{eq:interpolation} is a
nonzero element of $K$. It has order zero at every horizontal prime
$Q$ of $\bar Y$, so
\begin{equation}\label{eq:horizontal-interpolation}
 k\bigl(\ord_Q(\bar\pi^*s_i)-\ord_Q(\bar\pi^*s_1)\bigr)
        =(m_i-m_1)\ord_Q(\rho).
\end{equation}
The left side is bounded below independently of $i$, since each
$\bar\pi^*s_i$ is regular. As $m_i\to\infty$, the right side could
have this lower bound only if $\ord_Q(\rho)\ge0$. This proves
\textup{(i)}. If $S$ is a point, every prime is horizontal;
normality makes $\rho$ regular, and birational pushdown gives the
claimed section on $Z$. In this case the rank assertion also follows
from the section-ratio property. Indeed, the adjoint space in
\eqref{eq:induction-ranks} identifies with $H^0(Z,D+jkL)$:
the Jacobian divisor is effective and exceptional, so a rational
section allowed to have poles only there extends downstairs across
codimension two on smooth $Z$. The section $s_D\rho^j$ is nonzero,
and any two sections of $D+jkL$ have ratio in $K=\C$ by
\eqref{eq:systems}. Hence this space has dimension one for every
$j\ge0$, including $j=0$.

Assume $s>0$. We now verify \eqref{eq:induction-ranks} in every
nonnegative degree. The identity
\begin{equation}\label{eq:adjoint-identification}
 K_{Y/S}+\pi^*(J+jkL)
       =K_{Y/Z}+\pi^*(D+jkL)-f^*K_S
\end{equation}
shows which exceptional divisor is involved. Over the generic point
of $S$, a base canonical frame is immaterial. The Jacobian section,
$s_D$, and $\rho^j$ give a nonzero section of the right side.
To justify regularity, \textup{(i)} implies that $\rho$ is regular
on $\bar Y$ over a dense base open: remove the images of its finitely
many vertical polar divisors. Its pullback to $Y$ is then regular
there. Thus the adjoint space has dimension at least one, also for
$j=0$.

Take any two nonzero sections of that generic-fiber space. Choose a
rational frame of $K_S$ and regard them, by
\eqref{eq:adjoint-identification}, as rational sections of
$K_{Y/Z}+\pi^*(D+jkL)$ with no horizontal poles. Their finitely
many vertical polar divisors can be cleared by one effective Cartier
divisor $C$ on $S$: enclose their images, including any of codimension
at least two, in a Cartier divisor and increase its multiplicities so
that $f^*C$ dominates both polar divisors. Since $H_0$ is big, there
are an integer $a>0$ and a nonzero section of $aH_0-C$.
Multiplication by this section and by $s_C$ makes both rational
sections regular after the common base twist $af^*H_0$.

Now multiply both by the canonical section of $aR_0$. Formula
\eqref{eq:fixed-moving} produces two regular sections of
\[
                   K_{Y/Z}+\pi^*(D+jkL+aB_0)
\]
with the same ratio as the original pair. They descend to sections
of $D+jkL+aB_0$ on $Z$. In detail, $K_{Y/Z}$ is effective and
exceptional, and $\pi_*\OO_Y(K_{Y/Z})=\OO_Z$: a rational function
with poles only on exceptional divisors is regular away from a
codimension-two subset of smooth $Z$ and hence regular everywhere.
Apply this fact after trivializing the line bundle on $Z$.
The descended line belongs to \eqref{eq:systems}; the section-ratio
property therefore places their ratio in $K=\C(S)$.
Any two nonzero vectors in the generic-fiber section space are thus
proportional over its ground field. This proves dimension one for
all $j\ge0$, without discarding the zeroth degree.
\end{proof}

The same maximal system supplied both inputs for transfer. The
interpolation identity removed horizontal poles, while the big line
$H_0$ on the chosen base let us clear arbitrary vertical poles in
the rank calculation. No equidimensionality of the smooth resolution
was used: poles over higher-codimension base sets were cleared before
pushing down the exceptional Jacobian.

\subsection{Induction on dimension}

\begin{proof}[Proof of Theorem~\ref{thm:finite-volume}]
We induct on $d=\dim Z$. In dimension zero $Z$ is a point, every
line bundle is trivial, and the conclusion holds with $a=1,b=0$.
Assume $d>0$ and suppose that no $aL+bD$, with integers $a>0$ and
$b\ge0$, has a section.

The bundle $L$ is pseudoeffective because it has a semipositive metric.
The class of $L+D$ is therefore pseudoeffective. It is not big:
otherwise some positive multiple of $L+D$ would have a nonzero
section. Corollary~\ref{h:slope:face} gives a nonzero movable class
$\alpha$ with $L\cdot\alpha=D\cdot\alpha=0$ and
$\mu_{\min,\alpha}(T_Z)\ge0$. Apply Lemma~\ref{lem:determinant}
to obtain the fixed line $M$ and unbounded sections. Then apply
Lemma~\ref{lem:maximal-base}.

If its base is a point, the resulting section contradicts our
assumption. Otherwise Proposition~\ref{prop:transfer}, applied to
that prepared diagram and the two original metrics, supplies a line
$A$ and an effective integral divisor $D_S$ on the smaller smooth
projective base $S$. The line $A$ has a semipositive metric with
locally bounded weights, while $-K_S+D_S$ has a semipositive metric
with finite distinguished adjoint volume.

The no-forms condition also passes to $S$. A holomorphic form on
$S$ pulls back along the dominant rational map on its domain of
definition. Properness of $S$ extends the map at the generic point
of every divisor of $Z$, so that domain can be taken to have
codimension-two complement. The pulled-back form extends across that
complement because $\Omega_Z^q$ is locally free. Dominance in
characteristic zero makes pullback injective. Hence
$H^0(S,\Omega_S^q)=0$ for every $q>0$.

Since $\dim S<d$, induction produces integers $a>0,b\ge0$ and a
nonzero section of $aA+bD_S$. The lifting conclusion of
Proposition~\ref{prop:transfer} gives a nonzero section of
$akL+cD$ on $Z$ for some $c\ge0$. Its exponent of $L$ is positive,
so this contradicts the assumption and completes the induction.
\end{proof}

\section{Invariant sections and anticanonical nonvanishing}
\label{sec:torus}

The ordinary theorem is now available on every no-forms base. We apply
it to a rational torus quotient in order to obtain a section with the
invariance required by Section~\ref{sec:structure}. The quotient
provides adjoint rank one for a geometric reason: a rational function
constant on a dense orbit is a function on the base. Smoothness of the
anticanonical metric provides positivity of that invariant line.
We first establish the required smooth variant of transfer, then
verify its hypotheses on the quotient.

\subsection{The smooth invariant transfer}
\label{sec:invariant-transfer}

On a torus quotient it is the space of invariant adjoint sections,
rather than the full section space, that is one-dimensional.
We need a different positivity input for that line, while the divisor
normalization and boundary estimates remain those already proved.

\begin{proposition}[Smooth invariant variant]\label{prop:invariant-transfer}
Use the varieties, maps, line bundles, and section of
Proposition~\ref{prop:transfer}, with \(D=0\).
Let a compact torus \(T\) act holomorphically and continuously on
\(X,\bar Y,Y\), with compatible holomorphic linearizations of the
line bundles. Require \(r,\bar\pi,\bar f\) to be equivariant, with
trivial action on \(S\). Require \(\rho\) and the \(J\)-valued adjoint
section \(\eta_Y=\pi^*s_0\) to be invariant, and require \(h_J,h_L\)
to be smooth, semipositive, and \(T\)-invariant.

Replace the ordinary rank hypothesis \eqref{eq:rank} by
\[
 \dim_{\C(S)}
 H^0\!\left(Y_\eta,
 \left(K_{Y/S}+\pi^*(J+jkL)\right)|_{Y_\eta}\right)^T=1
 \qquad(j=0,1,2,\ldots).
\]
Then all the conclusions of Proposition~\ref{prop:transfer} hold,
and the section on \(X\) produced by its lifting is \(T\)-invariant.
\end{proposition}

\begin{proof}
The normalized line \(A\) is given the trivial action.
In \eqref{eq:minima} only base divisors are used, so dividing the
invariant \(\rho\) by their rational section gives an invariant \(e\).
The sections \(u_j\) in \eqref{eq:moments} are invariant as well.

On a dense smooth-fibration open where the adjoint direct image is
locally free and its formation commutes with base change, Berndtsson's theorem
\cite[Theorem~1.2]{Berndtsson2009} gives its smooth \(L^2\) metric
semipositive curvature. Its hypotheses hold because the total space
\(Y\) is projective, hence K\"ahler, and each twisting metric
\(\pi^*(h_Jh_L^{jk})\) is smooth and semipositive.

Average the holomorphic group action on this vector bundle over
the compact torus. This gives a holomorphic projection onto its
invariant subbundle, which has rank one by assumption.
The fiber integral is invariant under the group, so the projection
is orthogonal. The metric on the invariant line is therefore its
quotient metric and is semipositive. Since \(u_j\) is its nonzero
local generator, \(-\log V_j\) has the required psh representative.

This is the only positivity step that changes. The fixed-open
estimate \eqref{eq:common-open}, the zeroth-moment extension
\eqref{eq:base-volume-weight}, the high-moment envelope
\eqref{eq:envelope}, and both bounds through codimension two
now apply exactly as in the ordinary proof.
They give the two asserted base metrics.
Every base section is invariant because the base action is trivial.
The rational lifting multiplies its pullback by \(e^a\), so remains
invariant. With \(D=0\), the support condition permits only
exceptional poles, and the final extension on \(X\) is an invariant
holomorphic section.
\end{proof}

The distinction between the two propositions is substantive.
The ordinary argument allows a singular volume metric because the
whole adjoint space has rank one and the relative Bergman theorem
applies. The invariant argument uses smooth coefficients to obtain
a positive quotient metric on an orthogonal holomorphic summand.

\subsection{Natural invariant nonvanishing}

\begin{theorem}[Natural invariant nonvanishing]\label{thm:invariant}
Let $Z$ be smooth connected projective over $\C$, with
$H^0(Z,\Omega_Z^q)=0$ for every $q>0$. Suppose that
$L=-K_Z$ has a smooth Hermitian metric of semipositive curvature.
Let a compact connected real torus $T$ act continuously on $Z$ by
holomorphic automorphisms, and give $L$ the natural linearization
induced by the differential of the action. Then
\[
                  H^0(Z,mL)^T\ne0
                  \quad\text{for some integer }m>0.
\]
\end{theorem}

\begin{proof}
If $Z$ is a point the natural action on its canonical line is trivial,
so the constant section proves the claim. Otherwise replace $T$ by
its image in $\operatorname{Aut}(Z)$. The no-forms hypothesis and
Hodge symmetry give $H^1(Z,\OO_Z)=0$. The algebraic-torus argument
in the proof of Theorem~\ref{thm:structure} applies to this action:
connectedness preserves the first Chern class of a very ample line,
the vanishing of $H^1(\OO_Z)$ makes that line invariant up to
isomorphism, and its projective representation embeds $T$ in a
compact diagonal torus. Its Zariski closure $G$ is an algebraic torus
acting on $Z$, and natural sections are $T$-invariant exactly when
they are $G$-invariant.

Average the given metric over $T$ by averaging local weights, using
the natural linearization. Differences of metric weights are global
functions, so these averages patch to a smooth invariant metric; its
curvature remains semipositive. Take $J=L$ and $D=0$. The canonical
section $1$ of $K_Z+L=\OO_Z$ is invariant for the natural action,
and its metric pairing is a smooth positive volume form of finite
total mass on compact $Z$.

If the general $G$-orbit has dimension zero, connectedness of $G$
makes its action trivial: the orbit map is constant on a dense open
and hence on $Z$. Theorem~\ref{thm:finite-volume}, with $D=0$ and
$J=L$, then gives the required section. Suppose henceforth that the
general orbit has dimension $r>0$.

Take the Rosenlicht rational quotient
\cite{Rosenlicht1956,BellGhiocaReichstein2017}. Its function field is
$\C(Z)^G$, and a dense invariant open has geometric fibers equal
to orbits. Because $G$ is connected, its general orbit is geometrically
integral; thus the invariant field is relatively algebraically closed
in $\C(Z)$ and the quotient has connected general fiber. Choose its
prepared equivariant model
\[
      Y\xrightarrow{r_Y}\bar Y\xrightarrow{\bar\pi}Z,
          \qquad f:Y\to S,\quad\bar f:\bar Y\to S,
\]
with trivial base action. The notation $r_Y$ for the resolution avoids
confusion with the orbit dimension $r$. The generic fiber has dimension
$r$ and a geometrically dense $G$-orbit. After shrinking a base open,
$f$ is smooth there, its fibers have a dense orbit, and $\pi$ agrees
birationally with the original quotient construction.

We construct the relative section required for invariant transfer.
Choose $r$ elements of the Lie algebra of $G$ whose fundamental vector
fields $v_1,\ldots,v_r$ are generically independent. Such a choice is
possible because the general orbit has dimension $r$. The vector
fields are invariant, since $G$ is commutative, and tangent to the
fibers of $f$. On the smooth-fibration open their wedge is a regular,
generically nonzero section of $-K_{Y/S}$. Tensor it with the pullback
of a rational frame of $-K_S$, regular after shrinking the base open.
This gives an invariant anticanonical section on $Y$ over that open.
The determinant of $d\pi$ is the natural morphism
\[
                     -K_Y\longrightarrow\pi^*(-K_Z).
\]
Its image is a nonzero invariant rational section of $\pi^*L$,
regular over the same open.

This rational section also defines a section $\rho$ of
$\bar\pi^*L$ on $\bar Y$ with no horizontal poles. Indeed,
$r_Y$ is an isomorphism at the generic point of each prime divisor
of normal $\bar Y$. Every horizontal prime meets the base open on
which the constructed section is regular. The birational identification
of function fields therefore gives nonnegative order there.

The generic-fiber adjoint spaces are finite-dimensional
$\C(S)$-vector spaces carrying algebraic representations of the
split torus $G_{\C(S)}$. They decompose into character spaces.
A character of $G$ is trivial on the Zariski-dense subgroup $T$
exactly when it is the zero character. Consequently their
$T$-fixed and $G$-fixed subspaces agree, not merely those of sections
on $Z$ itself.

For each $j\ge0$, on the generic fiber the section
\[
             (\pi^*1/f^*ds)\rho^j
\]
is a nonzero invariant section of
$K_{Y/S}+\pi^*((j+1)L)$. The form $\pi^*1$ includes the effective
birational Jacobian, and $\rho$ has no horizontal pole, so this is
regular. Any two nonzero invariant sections of this line have an
invariant rational function as their ratio. The equivariant birational
maps and the defining property of the quotient give
\[
                    \C(Y)^G=\C(Z)^G=\C(S),
\]
so this ratio belongs to $\C(S)$. Their
space has dimension one over $\C(S)$, for every $j\ge0$ including
$j=0$.

Proposition~\ref{prop:invariant-transfer} now applies with $k=1$, with powers
subsequently taken there to clear the vertical minima. It supplies a
line $A$ and effective divisor $D_S$ on $S$, with a bounded
semipositive metric on $A$ and a finite-volume semipositive metric on
$-K_S+D_S$. Holomorphic forms on $S$ pull back injectively to $Z$,
as in the last part of Section~\ref{sec:induction}; hence $S$ has no
positive-degree holomorphic forms. Theorem~\ref{thm:finite-volume}
gives a nonzero section of $aA+bD_S$ with $a>0,b\ge0$.
Invariant lifting produces a nonzero invariant section of a positive
power of $L$ on $Z$. Since $D=0$, every permitted pole upstairs is
exceptional over $Z$ and disappears on its normal pushdown.
When $S$ is a point the same conclusion is the point-base case of
Proposition~\ref{prop:invariant-transfer}. This completes the proof.
\end{proof}

\subsection{Descent to the original variety}

\begin{proof}[Proof of Theorem~\ref{thm:headline}]
Use Theorem~\ref{thm:structure} to choose the finite \'etale cover
$\nu:S\to X$ and the universal product
$\widetilde X=\C^a\times F\times Z$ with its lattice deck action.
The compact factor $Z$ satisfies the hypotheses of
Theorem~\ref{thm:invariant}; let $s_Z\ne0$ be an invariant section
of $K_Z^{-m}$ for some $m>0$. If $Z$ is a point, take $m=1$ and
$s_Z=1$.

Using the product identification \eqref{eq:cover-canonical}, form
\[
     \eta_0^{-m}\otimes\eta_F^{-m}\otimes s_Z
                \in H^0(\widetilde X,K_{\widetilde X}^{-m}),
\]
where pullbacks by the product projections are understood. The first
two frames are invariant under the deck lattice and the last factor
is invariant under the closure of its action. The product is therefore
invariant and descends to a nonzero holomorphic section of $K_S^{-m}$.
It is algebraic by projectivity. Finally Lemma~\ref{lem:etale-norm}
produces a nonzero section of $K_X^{-m\deg\nu}$. The exponent is
positive, as required.
\end{proof}

\bibliographystyle{amsplain}
\bibliography{references}
\end{document}
